\section{Design}
Actuators were constructed using a combination of a folded polyimide sleeve and a internal  multilumen silicon tube with nitinol wires for actuation. Our methodology is inspired by similar "bellows-style" approaches to construct continuum devices (references) with the key difference in the use of Polyimide sheets (Kapton) forming the main body of the device.

This work focuses on the use of Kapton (thickness = ) due to its good robustness and high temperature capabilities.

\subsection{Construction of a Folded Protective Sleeve}
The polyimide film was folded into a []-style design as based on the previously published work [reference]. To achieve the folding a 2D drawing of the fold lines was drawn in Fusion360 and printed onto paper. This paper was taped to the table and a sheet of Kapton was taped over the top, the Kapton is clear so the printed lines are visible through. The lines were then traced onto the Kapton with a ruler and a ballpoint pen. The ballpoint pen marked the desired fold lines and also partially indented on the lines. The sheet of kapton was removed from the table and the outer edge was cut with scissors. The Kapton sheet was then folded along the lines as shown in the figure [ref figure]. The folds on the sheet were puished back out and the kapton was rolled into a cylinder. This was taped along the overlapped edge to form a temporary seal while maintaining care to line up the folds. The tape was removed in small sections and JBweld HighHeat was applied until the whole adhesion region. This was threaded onto a plastic tube with a diameter approx. 2mm smaller than the kapton sleeve. The sleeve was taped to the tube and the adhesion region was taped over the top to clamp the glued region. This was left for 24 hours to cure. The glueing and the curing process was performed in a Fume cupboard [brand/details] due to the risk of respiratory irritation through the curing process. After curing, a strip of Kapton tape was placed over the edge on the adhesion region. This improved the flexiblity of the region. 


For the proposed design, the film selected was Kapton\textsuperscript{\textregistered} ( Supplier, thickness). 

\subsubsection{Manual folding versus using a plotter}
A commercially available plotter (Cricut\textsuperscript{\textregistered}) was used with a scribe tool and a cutter to achieve a rapid manufacturing of a folded sleeve. This was attempted with paper initally for cheap prototyping and then transferring to the polymer film. 

The paper was easy however the polyimide did not take up the scores so easily. We demonstrated this difference by using the plotter to score a 5cm line on piece of 80gsm paper (roughly 0.25mm) and two pieces of Kapton (thicknesses). These were attached to the jaws of a vice and compressed to form a crease on the scored line. On the paper attempt, after the material was allowed to relax, a defined crease was left in the material. In comparison, the Kapton did not take up the crease and when relaxed, remianed in a curve.

\subsection{Complete Assembly}
The complete assembly was formed with the polyimide sleeve, the silicon continuum device and 3D printed parts. The construction method is as follows: a sheet of polyimide is marked with the desired folding pattern (as shown in figure), the sheet is folded by hand before being shaped into a tube and glued along an overlapping edge. The folds are encouraged into place and one end is sealed. This is the assembled with 3D printed parts to form the actuator.

The silicon actuator contains 7 holes where three lengths of nitinol wire were threaded through and fastened with two grub screws each in 3D printed handles. This leaves a separate tube through the centre of the silicon tube that can be used for targeted delivery of liquid or components such as naneye cameras for observation.

